\SourceLeaf{372}%
It is hard to understand how he could have deceived himself in this way. He even raises the objection himself that if the pieces \(Sp,pp'\) and so on become smaller and smaller, their sum might still never grow beyond a certain limit, although their number increases without end. He thinks he can answer this objection by saying that even as those pieces become smaller and smaller, they always remain greater than \textit{an assignable magnitude}. For those pieces are legs of right triangles, and hence are always greater than the difference between the hypotenuse and the other leg. It almost seems that the author was misled by a grammatical ambiguity, namely the two meanings of the article \textit{an} in an assignable magnitude. The author's inference would be valid only if it could be shown that the pieces \(Sp,pp'\) and so on always remain greater. than \textit{a fixed} assignable magnitude, for example the difference between the hypotenuse \(pT\) and the leg \(ST\). All that can be proved, however, is that each piece always remains greater than an assignable magnitude that is different for each piece. Thus \(Sp\) is greater than the difference between \(pT\) and \(ST\), \(pp'\) is greater than the difference between \(qp'\) and \(qp\), and so on. With this the proof loses all its force.

By the same method he thought would establish his proof, he could also prove that in a plane triangle \(ABC\) in which \(B\) is a right angle, \(C\) cannot be acute. He would only have to drop from \(B\) a perpendicular \(BD\) to the hypotenuse \(AC\), then drop the perpendicular \(DE\) to \(AB\), and go on without end dropping the perpendiculars \(EF,FG,GH\) and so on alternately to \(AC\) and \(AB\). The pieces \(CD,DF,FH\) and so on are always greater than the assignable difference betweenhypotenuse and one leg of \textit{the particular} right triangle in which each in turn is the other leg. Yet their sum clearly never exhausts the hypotenuse \(AC\), however large their number is taken to be.

We would almost have to regret dwelling so long on such familiar and easy matters, were it not for the more than usual attention this work drew from the way it was announced in the public press even before it appeared. The author, incidentally, seems genuinely concerned with the truth. We will therefore add only that the author subsequently tries, by a very similar and therefore equally ineffective argument, to prove %
